%==================================================================
%  Section 1: Introduction.
%=====================================================================
\section{Introduction}
\label{sec:intro}
Supervised fine-tuning (SFT) is a central stage in modern open-weight
LLM post-training recipes such as Llama, Qwen, and
Mistral~\cite{dubey2024llama3,qwen2025technical,jiang2023mistral}.
Although these recipes differ in scale, supervision, and alignment
objectives, they share a common dependency: the SFT data determines
the first task-facing adaptation of the base model and shapes the
model state on which later preference optimization or reinforcement
learning operates. High-quality instruction data is therefore a
primary determinant of effective adaptation, making its selection a
practical bottleneck for efficient and reliable post-training. Yet
quality is not a single scalar: a candidate must both provide
credible supervision and offer learning value to the model at its
current training state.

\paragraph{Two axes of instruction-data quality.}
These requirements define two related but non-substitutable axes.
\emph{Reliability} asks whether a response provides credible
supervision for its instruction, whereas \emph{usefulness} asks
whether learning from that credible pair benefits the current model.
A reliable example may already be mastered or redundant, and its
learning value may change as other capabilities are acquired.
Reliability is therefore a pair-level admissibility question that can
be assessed once before fine-tuning; usefulness is model-dependent
and must be re-evaluated as training progresses. This distinction
implies an ordered selection process: first screen or suppress
unreliable supervision, then dynamically rank the remaining
candidates by their current learning value.

\paragraph{From static scoring to dynamic usefulness.}
A one-time assessment is appropriate for reliability but insufficient
for model-dependent usefulness. Most instruction-data selection
methods score and rank candidates once before fine-tuning.
Quality-based methods use external judges or learned quality
estimators; target-conditioned methods estimate influence relative
to a validation or few-shot target set; and representation-based
methods project candidate activations onto target-derived hidden-state
directions. A more recent line of dynamic data selection instead
re-estimates candidate utility as the model state changes during
training~\cite{kung2023activeit, sun2025dots}. This line highlights a
key fact: an example's usefulness is not a static property of the
data alone, but depends on the current model, training stage, and
optimisation regime. Once unreliable supervision has been screened,
dynamic selection is therefore the critical mechanism for identifying
which remaining examples offer the greatest learning value at the
current stage. We focus on supervised instruction tuning, where
candidates are instruction--response pairs and no rollout group,
verifier reward, or target validation set is assumed.

\paragraph{Three gaps in reliable training-adaptive selection.}
Despite this progress, three gaps remain. First, existing dynamic
pipelines generally rank the full candidate pool without first
establishing whether its supervision is trustworthy. A static
reliability gate serves an orthogonal role: it asks whether a pair is
admissible before the dynamic selector asks how useful that pair is
to the current model. It is not another model-dependent difficulty or
uncertainty metric, but a new pre-selection stage that establishes a
credible candidate pool for dynamic ranking.

Second, within that credible pool, structural hidden-state
selection~\cite{nait2026} and dynamic
selection~\cite{kung2023activeit, sun2025dots, dataagent}
have largely evolved separately. Hidden-state directions are
typically fixed before training~\cite{nait2026}, while dynamic
rankings rely on scalar utilities such as difficulty or uncertainty.
Consequently, a selector can track scalar utility without also using
a representation-geometric signal recomputed from the current model
state.

Third, scalar utilities can be sensitive in heterogeneous multi-task
instruction pools. In composite-reward selectors, the variance-ratio
weight that mixes difficulty and uncertainty can reflect between-task
scale differences as well as within-task informativeness. This
confounding calls for a role-separated fusion in which the
representation signal enters as an explicitly bounded, interpretable,
and independently ablatable modulation.

The absence of the first stage is especially consequential in
low-quality instruction pools, where mismatched pairs, truncated
responses, and subtly wrong answers are
common~\cite{honovich2023unnatural,alpagasus}. The dynamic
utility signals above presuppose that the pair itself can be trusted,
yet difficulty- or uncertainty-driven selection can actively
\emph{prefer} corruption because it can raise loss and entropy in the
same direction as a genuinely hard clean pair. This is the
instruction-tuning analogue of the hard-versus-mislabeled ambiguity
long documented for
classification~\cite{pleiss2020aum,swayamdipta2020cartography,rholoss2022}:
hardness scores conflate ``hard because informative'' with ``hard
because broken.'' A reliability gate addresses this prior question
on an orthogonal axis; it does not replace dynamic selection, but
makes dynamic selection meaningful by restricting utility
comparisons to examples with credible supervision.

\paragraph{Our approach: reliability-gated, state-conditioned
selection.}
We propose \textsc{TAG} (\textbf{T}raining-\textbf{A}daptive data
selection with a reliability \textbf{G}ate). \textsc{TAG} brings the
two quality axes together in a three-view selector: a static
reliability estimate and two training-adaptive utility signals. The
gate assigns each instruction--response pair a continuous reliability
weight \(G_i\in[0,1]\) at the base checkpoint and caches it. At each
refresh step~\(t\), the current model supplies a
difficulty--uncertainty reward \(\mathcal R_i^{(t)}\) and a
state-conditioned representation-anchor signal, giving the score
\begin{equation}
s_i^{(t)}
=
\underbrace{G_i}_{\substack{\text{reliability gate}\\ \text{(static, once)}}}
\;\cdot\;
\underbrace{\mathcal{R}_i^{(t)}}_{\text{dynamic reward}}
\;\cdot\;
\underbrace{\bigl(1+\lambda\,\widetilde a_i^{(t)}\bigr)}_{\text{bounded anchor modulation}}.
\label{eq:tag-score}
\end{equation}
At each refresh, \(\widetilde a_i^{(t)}\in[0,1]\) is the
candidate's normalized projection onto per-layer principal
representation-displacement directions extracted from a current
probe subset. We use these sign-oriented directions as a
state-conditioned representation anchor. The parameter
\(\lambda\ge0\) controls the modulation strength, so the anchor
factor is confined to \([1,1+\lambda]\) and can be ablated by
setting \(\lambda=0\). The multiplicative gate is not a hard binary
filter: for \(G_i>0\) it continuously attenuates each candidate
according to its estimated reliability, while its zero anchor makes
\(G_i=0\) an exact veto that no amount of dynamic utility can
compensate; the two dynamic terms determine a credible pair's
usefulness to the evolving model.

\paragraph{Design rationale.}
Three choices support this construction. First, the reliability gate
acts as a detector of instruction--response dependency rather than a
classifier of corruption types: an instruction-side counterfactual
asks how strongly the response depends on its paired instruction and
maps that evidence to a continuous reliability gate. Second,
because the representation anchor is recomputed during training, its
local directional stability matters. Under optimiser-scale covariance
drift and positive cross-time spectral separation, a Davis--Kahan
analysis gives a conditional \(O(\eta_t)\) one-step bound on the
sign-oriented principal direction. This does not guarantee stability
of candidate projections or of the full ranking; the bounded factor
instead limits the anchor's direct modulation of the reward. Third,
the dynamic score is
analytically motivated: a total-variance decomposition shows that a
difficulty--uncertainty mixture can entangle within-task
informativeness with between-task scale differences, motivating the
addition of a distinct and explicitly bounded representation signal.

\paragraph{Contributions.}
We make four contributions:
\begin{itemize}
  \item We introduce \textsc{TAG}, a training-adaptive selection
  framework that combines a continuous, static reliability gate
  with dynamic estimates of model-dependent usefulness.
  \item We develop an instruction-side counterfactual detector that
  aggregates response-level, span-level, and completeness evidence
  into a reliability gate \(G_i\in[0,1]\), computed once and cached.
  \item We provide analytical support for the dynamic selector: a
  total-variance decomposition motivates the role-separated reward
  design, and a conditional local stability result characterizes the
  recomputed principal direction, whose normalized candidate projection
  provides a bounded modulation of that reward.
  \item We evaluate \textsc{TAG} on instruction pools with typed,
  manifest-verified corruptions (mismatched, noisy, truncated,
  wrong-answer, and duplicated), measuring both corruption detection
  (Dirty@\(K\)/AUPRC) and end-to-end SFT quality against static and
  dynamic selection baselines; \STATUS{results pending}.
\end{itemize}




